\section{Conclusion}\label{sec:conclusion}

Mature theories keep introducing small sets of variables that bound feasibility, add up along protocols, or act as exchange rates between competing aims. Within the Six Birds closure calculus we have treated these variables as currencies, defined by the structural role they play relative to a layer.

The central result is conditional and exact. Suppose a lower-layer ledger must stay within a budget for higher-layer objects to be maintained. Then the ledger is a feasibility constraint for the higher layer. If that layer closes by maximum entropy, the constraint's Gibbs multiplier is a nonnegative shadow price. Provided some weighted row has nonconstant costs, that price is positive exactly when the budget lies strictly between the cheapest feasible cost and the unconstrained cost, zero at or above the unconstrained cost, and divergent as the budget approaches the cheapest feasible cost; with constant costs the price is not identifiable. Path-reversal KL, and its reversal-mixture regularization on finite samples, cannot be inflated by coarse observation. And with stationary prototypes, a fixed horizon, and positive within-block mass in every row, the uniform bound on packaging defects tends to zero as the price of leaving a block grows.

A finite Markov laboratory showed each step at work. It also marked where the evidence is a theorem, where it is a property of one construction, and where an apparent pattern does not hold in general. The broader picture follows. Wherever stable description must be maintained on a limited resource, budgets appear, and where budgets bind in a maximum-entropy closure, prices appear. Persistence alone does not supply the budget; given a feasible budget and such a closure, the price follows.
